\Needspace{12\baselineskip}
\section{References}

\begingroup
\setlength{\parskip}{0pt}
\titlespacing*{\subsection}{0pt}{1.0ex}{0.3ex}

\subsection*{Scripture}

{\fontsize{8.6}{9.6}\selectfont
\begin{itemize}[itemsep=0pt,parsep=0pt,topsep=0.05em]
\item Ecclesiasticus 7:40. 2~Machabees 12:43--46. Matthew 5:25--26; 7:21--23; 11:22--24; 12:36--37; 13:36--43, 47--50; 16:27; 24:29--36; 25:1--46. Luke 12:2--5; 16:19--31; 23:42--43. John 3:17--21; 5:22--30; 11:24--26; 12:47--50. Acts 17:30--31. Romans 2:5--16; 14:10--12. 1~Corinthians 3:11--17; 4:3--5; 15:20--28, 51--55. 2~Corinthians 5:1--10. Philippians 1:21--24. 1~Thessalonians 4:13--17. 2~Timothy 4:1, 6--8. Hebrews 9:27--28. Apocalypse 14:13; 20:11--15.
\item All quoted from the public-domain Douay--Rheims translation in Bishop Challoner's revision, in the tracked per-book verse text of Project Gutenberg eBook~1581 registered in this repository's reusable source library; wording checked 2026-07-25, except Ecclesiasticus 7:40, located and checked 2026-07-26. Chapter and verse follow the Douay--Rheims's own numbering, which follows the Clementine Vulgate. The registered divergence table for that text marks Matthew 24:36 and John 5:22 as differing from the settled American edition of 1899; both divergences are stated where the verses are quoted. No translation was composed for this article and no other English version is quoted in it.
\end{itemize}}

\subsection*{Patristic witness}

{\fontsize{8.6}{9.6}\selectfont
\begin{itemize}[itemsep=0pt,parsep=0pt,topsep=0.05em]
\item Augustine of Hippo, \emph{De civitate Dei} XX.1 (why the last day is properly the day of judgment; ``day'' used for ``time''), XX.2 (the hiddenness of God's present judgments), XX.14 (the books of Apocalypse~20; the book of life as a divine power recalling every man's works), XXI.13 (against those who make all punishment after death purgatorial), XXI.26 (1~Corinthians~3 read of the fire of tribulation in this life; the concession about a purgative fire in the interval). English: trans.\ Marcus Dods, \emph{The City of God}, vol.~II, in \emph{The Works of Aurelius Augustine, A New Translation} (Edinburgh: T.\,\&\,T. Clark, 1871); public domain; wording checked 2026-07-25 in this repository's tracked transcription of the 1871 printing, with passages registered for XX.1 (artifact lines 13879--13938), XX.14 (14884--14951), and XXI.26 (18318--18475).
\item Augustine of Hippo, \emph{De civitate Dei} XXI.26, Latin (\latin{non redarguo, quia forsitan uerum est}); ed.\ Emanuel Hoffmann, \emph{Corpus Scriptorum Ecclesiasticorum Latinorum} 40.2 (Vienna, 1900); public domain; checked 2026-07-25 in this repository's tracked TEI text of that edition. The edition prints \latin{u} for consonantal \latin{v}.
\item Augustine, \emph{Sermo} 18, 4 (PL 38, 130--131), on the least brethren as Christ's members, cited only as quoted at \emph{Catechism of the Catholic Church} 1039; the sermon was not read at its own locus.
\end{itemize}}

\subsection*{Thomistic framework}

{\fontsize{8.6}{9.6}\selectfont
\begin{itemize}[itemsep=0pt,parsep=0pt,topsep=0.05em]
\item Thomas Aquinas, \emph{Summa theologiae} III, q.~59, aa.~1--6, especially a.~5 with its three replies (whether a general judgment remains after the judgment that takes place in the present time; the five ways a finished life continues; the unchangeable state of the separated soul; the changeable body). English from the public-domain translation of the Fathers of the English Dominican Province at \sourceurl{https://www.newadvent.org/summa/4059.htm}{New Advent}; Latin at \sourceurl{https://www.corpusthomisticum.org/sth4053.html}{Corpus Thomisticum}; both checked 2026-07-25 and registered as exact artifacts of those dated web states.
\item Thomas Aquinas, \emph{Scriptum super libros Sententiarum} IV, d.~47, q.~1, a.~1, qc.~1, corpus and replies 1--3 (the twofold divine operation and the twofold judgment; \latin{unum singulare, quod de eo fiet post mortem}; the increase of reward after the judgment). Latin at \sourceurl{https://www.corpusthomisticum.org/snp4047.html}{Corpus Thomisticum}, checked 2026-07-25 and registered as an exact artifact of that dated web state.
\item \emph{Supplementum Tertiae Partis Summae theologiae} q.~69, aa.~1--2 (the places that receive souls; immediate conveyance unless held by a debt); q.~87, aa.~1--3 (each man's conscience as the book; the manifestation of merits to all; the rapidity of the survey); q.~88, aa.~1--4 (the general judgment; the mental conduct of the debate; the reserved time; the place); Appendix~I, q.~2, aa.~1--6 (the pains of purgatory); Appendix~II, q.~1, aa.~1--2 (whether there is a purgatory; where it is). The Supplement is not a work of Aquinas's composition: the Leonine editors state that it was transcribed about the middle of the fourteenth century from a codex of the commentary on the fourth book of the Sentences and is to be held a second witness to that commentary. English at \sourceurl{https://www.newadvent.org/summa/5088.htm}{New Advent} (with the pages for qq.~69 and~87 and the two appendices); Latin at the Internet Archive page images of \emph{Sancti Thomae Aquinatis Opera omnia}, tomus duodecimus (Rome: Typographia Polyglotta, 1906), leaf~412 (the \emph{Praefatio in Supplementum}) and leaf~664 (printed p.~205, q.~88 aa.~1--3), with the volume's uncorrected machine transcription used only as a finding aid; all checked 2026-07-25 and registered.
\end{itemize}}

\subsection*{Defined and authoritative Catholic teaching}

{\fontsize{8.6}{9.6}\selectfont
\begin{itemize}[itemsep=0pt,parsep=0pt,topsep=0.05em]
\item Fourth Lateran Council (1215), constitution~1 \emph{Firmiter credimus}, at DS 801 (the coming at the end of the age; rising \latin{cum suis propriis corporibus, quae nunc gestant}). Second Council of Lyons (1274), profession of faith of Michael Palaeologus, at DS 856--859 (purgatorial punishments and suffrages; \latin{mox} to heaven and to hell; \latin{nihilominus in die iudicii}). John~XXII, \emph{Ne super his} (3 December 1334), at DS 990--991. Benedict~XII, \emph{Benedictus Deus} (29 January 1336), at DS 1000--1002 (the intuitive and facial vision before the resumption of bodies and the general judgment; immediate descent for those dying in actual mortal sin; the retained day of judgment, quoting 2~Corinthians 5:10). Council of Florence, \emph{Laetentur caeli} (6 July 1439), at DS 1304--1306. Latin of all five read 2026-07-25 in the dated web presentation of the Latin \emph{Enchiridion symbolorum} at \sourceurl{https://patristica.net/denzinger/enchiridion-symbolorum.html}{patristica.net}, registered as an exact artifact of that state. The Enchiridion is a compilation and a finding aid to its underlying documents; no printed Denzinger edition and no conciliar manuscript was collated. English renderings in this article are working glosses except where the Catechism's own English is quoted.
\item Council of Trent, session~VI (13 January 1547), canons~30, 32, and~33 on justification; session~XXV (3--4 December 1563), \emph{Decree concerning Purgatory}. English: trans.\ J.~Waterworth, \emph{The Canons and Decrees of the Sacred and Oecumenical Council of Trent} (London 1848), in an undated Burns and Oates reprint; page images of printed pp.~48, 49, 232, and 233 inspected 2026-07-25 and registered. Latin: \emph{Canones et decreta sacrosancti oecumenici Concilii Tridentini}, editio stereotypa undecima (Leipzig: Tauchnitz, 1887), reprinting the Roman edition of 1834; page images of printed pp.~38 and~173 inspected 2026-07-25 and registered. The canons are DS 1580 and the decree DS 1820.
\item \emph{Catechism of the Catholic Church} 678--679 (Christ's judgment of the living and the dead), 997--1001 (the resurrection of the body), 1021--1022 (the particular judgment), 1030--1032 (the final purification), 1038--1041 (the Last Judgment), 1861 (judgment of persons entrusted to God). Official Vatican English, archived IntraText presentation, pages checked 2026-07-25 at \sourceurl{https://www.vatican.va/archive/ENG0015/__P2L.HTM}{vatican.va (1021--1022)} and the adjacent pages for the other paragraphs, all registered as exact artifacts of that delivery state.
\item Sacred Congregation for the Doctrine of the Faith, \emph{Recentiores episcoporum Synodi}, letter on certain questions concerning eschatology, 17~May 1979, points~1--7 and the pastoral section. Holy See English text at \sourceurl{https://www.vatican.va/roman_curia/congregations/cfaith/documents/rc_con_cfaith_doc_19790517_escatologia_en.html}{vatican.va}, checked 2026-07-25 and registered. An approved but non-definitive curial act; the Latin of \emph{Acta Apostolicae Sedis} 71 (1979) was not consulted.
\end{itemize}}

\endgroup
